% Expanded bibliography and source map for the Metatime monograph.
% Version v10.8, 2026-07-29.
% The grouping distinguishes external scientific context from original TIR claims.

\renewcommand{\bibname}{Bibliography and Source Map}
\begin{thebibliography}{999}
\small
\sloppy

\item[]\textit{Editorial note.} The references below document the established
mathematical, physical, experimental, and methodological context used throughout
this monograph. Inclusion of a source does not imply that the source endorses the
Metatime/TIR construction, nor that the original claims of this monograph follow
from the cited literature. Numerical particle data should be read against the
current Particle Data Group listings and their published updates.

\item[]\textbf{A. Reference data, precision measurements, and major experiments}

\bibitem{PDG2024}
S.~Navas \emph{et al.} (Particle Data Group),
``Review of Particle Physics,''
Phys.\ Rev.\ D \textbf{110}, 030001 (2024), with the 2025 online update.

\bibitem{CODATA2022}
E.~Tiesinga, P.~J.~Mohr, D.~B.~Newell, and B.~N.~Taylor,
``The 2022 CODATA Recommended Values of the Fundamental Physical Constants,''
NIST Standard Reference Database 121, Web Version 9.0 (2024),
\url{https://physics.nist.gov/constants}.

\bibitem{Planck2018}
N.~Aghanim \emph{et al.} (Planck Collaboration),
``Planck 2018 results. VI. Cosmological parameters,''
Astron.\ Astrophys.\ \textbf{641}, A6 (2020).

\bibitem{DESI2024}
A.~G.~Adame \emph{et al.} (DESI Collaboration),
``DESI 2024 VI: Cosmological Constraints from the Measurements of Baryon
Acoustic Oscillations,'' arXiv:2404.03002 (2024).

\bibitem{ATLAS2012Higgs}
G.~Aad \emph{et al.} (ATLAS Collaboration),
``Observation of a new particle in the search for the Standard Model Higgs
boson with the ATLAS detector at the LHC,''
Phys.\ Lett.\ B \textbf{716}, 1--29 (2012).

\bibitem{CMS2012Higgs}
S.~Chatrchyan \emph{et al.} (CMS Collaboration),
``Observation of a new boson at a mass of 125 GeV with the CMS experiment at
the LHC,'' Phys.\ Lett.\ B \textbf{716}, 30--61 (2012).

\bibitem{ATLAS2022Charm}
G.~Aad \emph{et al.} (ATLAS Collaboration),
``Direct constraint on the Higgs--charm coupling from a search for Higgs boson
decays into charm quarks with the ATLAS detector,''
Eur.\ Phys.\ J.\ C \textbf{82}, 717 (2022).

\bibitem{SuperK1998}
Y.~Fukuda \emph{et al.} (Super-Kamiokande Collaboration),
``Evidence for oscillation of atmospheric neutrinos,''
Phys.\ Rev.\ Lett.\ \textbf{81}, 1562--1567 (1998).

\bibitem{SNO2002}
Q.~R.~Ahmad \emph{et al.} (SNO Collaboration),
``Direct evidence for neutrino flavor transformation from neutral-current
interactions in the Sudbury Neutrino Observatory,''
Phys.\ Rev.\ Lett.\ \textbf{89}, 011301 (2002).

\bibitem{KamLAND2003}
K.~Eguchi \emph{et al.} (KamLAND Collaboration),
``First results from KamLAND: Evidence for reactor antineutrino disappearance,''
Phys.\ Rev.\ Lett.\ \textbf{90}, 021802 (2003).

\bibitem{DayaBay2012}
F.~P.~An \emph{et al.} (Daya Bay Collaboration),
``Observation of electron-antineutrino disappearance at Daya Bay,''
Phys.\ Rev.\ Lett.\ \textbf{108}, 171803 (2012).

\bibitem{Abel2020nEDM}
C.~Abel \emph{et al.},
``Measurement of the permanent electric dipole moment of the neutron,''
Phys.\ Rev.\ Lett.\ \textbf{124}, 081803 (2020).

\bibitem{MuonG2Final2025}
D.~P.~Aguillard \emph{et al.} (Muon $g-2$ Collaboration),
``Measurement of the positive muon anomalous magnetic moment to 127 ppb,''
Phys.\ Rev.\ Lett.\ \textbf{135} (2025), DOI: 10.1103/7clf-sm2v.

\bibitem{Riess1998}
A.~G.~Riess \emph{et al.},
``Observational evidence from supernovae for an accelerating universe and a
cosmological constant,'' Astron.\ J.\ \textbf{116}, 1009--1038 (1998).

\bibitem{Perlmutter1999}
S.~Perlmutter \emph{et al.},
``Measurements of $\Omega$ and $\Lambda$ from 42 high-redshift supernovae,''
Astrophys.\ J.\ \textbf{517}, 565--586 (1999).

\item[]\textbf{B. Quantum mechanics, quantum field theory, and the Standard Model}

\bibitem{Noether1918}
E.~Noether,
``Invariante Variationsprobleme,''
Nachr.\ Ges.\ Wiss.\ G\"ottingen, Math.-Phys.\ Kl. 235--257 (1918).

\bibitem{Dirac1928}
P.~A.~M.~Dirac,
``The quantum theory of the electron,''
Proc.\ R.\ Soc.\ Lond.\ A \textbf{117}, 610--624 (1928); \textbf{118},
351--361 (1928).

\bibitem{Weyl1929}
H.~Weyl,
``Elektron und Gravitation. I,''
Z.\ Phys.\ \textbf{56}, 330--352 (1929).

\bibitem{Tomonaga1946}
S.~Tomonaga,
``On a relativistically invariant formulation of the quantum theory of wave
fields,'' Prog.\ Theor.\ Phys.\ \textbf{1}, 27--42 (1946).

\bibitem{Schwinger1948}
J.~Schwinger,
``Quantum electrodynamics. I. A covariant formulation,''
Phys.\ Rev.\ \textbf{74}, 1439--1461 (1948).

\bibitem{Feynman1949}
R.~P.~Feynman,
``Space-time approach to quantum electrodynamics,''
Phys.\ Rev.\ \textbf{76}, 769--789 (1949).

\bibitem{Dyson1949}
F.~J.~Dyson,
``The radiation theories of Tomonaga, Schwinger, and Feynman,''
Phys.\ Rev.\ \textbf{75}, 486--502 (1949).

\bibitem{YangMills1954}
C.~N.~Yang and R.~L.~Mills,
``Conservation of isotopic spin and isotopic gauge invariance,''
Phys.\ Rev.\ \textbf{96}, 191--195 (1954).

\bibitem{Glashow1961}
S.~L.~Glashow,
``Partial-symmetries of weak interactions,''
Nucl.\ Phys.\ \textbf{22}, 579--588 (1961).

\bibitem{Weinberg1967}
S.~Weinberg,
``A model of leptons,'' Phys.\ Rev.\ Lett.\ \textbf{19}, 1264--1266 (1967).

\bibitem{Salam1968}
A.~Salam,
``Weak and electromagnetic interactions,'' in
\emph{Elementary Particle Theory}, edited by N.~Svartholm
(Almqvist and Wiksell, Stockholm, 1968), pp.~367--377.

\bibitem{EnglertBrout1964}
F.~Englert and R.~Brout,
``Broken symmetry and the mass of gauge vector mesons,''
Phys.\ Rev.\ Lett.\ \textbf{13}, 321--323 (1964).

\bibitem{Higgs1964}
P.~W.~Higgs,
``Broken symmetries and the masses of gauge bosons,''
Phys.\ Rev.\ Lett.\ \textbf{13}, 508--509 (1964).

\bibitem{GHK1964}
G.~S.~Guralnik, C.~R.~Hagen, and T.~W.~B.~Kibble,
``Global conservation laws and massless particles,''
Phys.\ Rev.\ Lett.\ \textbf{13}, 585--587 (1964).

\bibitem{tHooftVeltman1972}
G.~'t~Hooft and M.~Veltman,
``Regularization and renormalization of gauge fields,''
Nucl.\ Phys.\ B \textbf{44}, 189--213 (1972).

\bibitem{GrossWilczek1973}
D.~J.~Gross and F.~Wilczek,
``Ultraviolet behavior of non-Abelian gauge theories,''
Phys.\ Rev.\ Lett.\ \textbf{30}, 1343--1346 (1973).

\bibitem{Politzer1973}
H.~D.~Politzer,
``Reliable perturbative results for strong interactions?''
Phys.\ Rev.\ Lett.\ \textbf{30}, 1346--1349 (1973).

\bibitem{Wilson1974}
K.~G.~Wilson,
``Confinement of quarks,'' Phys.\ Rev.\ D \textbf{10}, 2445--2459 (1974).

\bibitem{PeskinSchroeder1995}
M.~E.~Peskin and D.~V.~Schroeder,
\emph{An Introduction to Quantum Field Theory}
(Addison-Wesley, Reading, 1995).

\bibitem{WeinbergQFT1995}
S.~Weinberg,
\emph{The Quantum Theory of Fields, Vols. I--III}
(Cambridge University Press, Cambridge, 1995--2000).

\item[]\textbf{C. Flavor, hadrons, neutrinos, and symmetry classification}

\bibitem{gellmann1961}
M.~Gell-Mann,
``The eightfold way: A theory of strong interaction symmetry,''
Caltech Synchrotron Laboratory Report CTSL-20 (1961).

\bibitem{Neeman1961}
Y.~Ne'eman,
``Derivation of strong interactions from a gauge invariance,''
Nucl.\ Phys.\ \textbf{26}, 222--229 (1961).

\bibitem{GellMann1962}
M.~Gell-Mann,
``Symmetries of baryons and mesons,''
Phys.\ Rev.\ \textbf{125}, 1067--1084 (1962).

\bibitem{okubo1962}
S.~Okubo,
``Note on unitary symmetry in strong interactions,''
Prog.\ Theor.\ Phys.\ \textbf{27}, 949--966 (1962).

\bibitem{Cabibbo1963}
N.~Cabibbo,
``Unitary symmetry and leptonic decays,''
Phys.\ Rev.\ Lett.\ \textbf{10}, 531--533 (1963).

\bibitem{GellMann1964}
M.~Gell-Mann,
``A schematic model of baryons and mesons,''
Phys.\ Lett.\ \textbf{8}, 214--215 (1964).

\bibitem{Zweig1964}
G.~Zweig,
``An SU(3) model for strong interaction symmetry and its breaking,''
CERN Reports TH-401 and TH-412 (1964).

\bibitem{kobayashi1973}
M.~Kobayashi and T.~Maskawa,
``CP-violation in the renormalizable theory of weak interaction,''
Prog.\ Theor.\ Phys.\ \textbf{49}, 652--657 (1973).

\bibitem{Wolfenstein1983}
L.~Wolfenstein,
``Parametrization of the Kobayashi--Maskawa matrix,''
Phys.\ Rev.\ Lett.\ \textbf{51}, 1945--1947 (1983).

\bibitem{Jarlskog1985}
C.~Jarlskog,
``Commutator of the quark mass matrices in the Standard Electroweak Model and
a measure of maximal CP nonconservation,''
Phys.\ Rev.\ Lett.\ \textbf{55}, 1039--1042 (1985).

\bibitem{Maki1962}
Z.~Maki, M.~Nakagawa, and S.~Sakata,
``Remarks on the unified model of elementary particles,''
Prog.\ Theor.\ Phys.\ \textbf{28}, 870--880 (1962).

\bibitem{Pontecorvo1957}
B.~Pontecorvo,
``Mesonium and antimesonium,''
Zh.\ Eksp.\ Teor.\ Fiz.\ \textbf{33}, 549--551 (1957)
[Sov.\ Phys.\ JETP \textbf{6}, 429 (1958)].

\bibitem{Wolfenstein1978}
L.~Wolfenstein,
``Neutrino oscillations in matter,''
Phys.\ Rev.\ D \textbf{17}, 2369--2374 (1978).

\bibitem{MikheyevSmirnov1985}
S.~P.~Mikheyev and A.~Y.~Smirnov,
``Resonance amplification of oscillations in matter and spectroscopy of solar
neutrinos,'' Sov.\ J.\ Nucl.\ Phys.\ \textbf{42}, 913--917 (1985).

\item[]\textbf{D. Geometric phase, differential geometry, topology, and holonomy}

\bibitem{Pancharatnam1956}
S.~Pancharatnam,
``Generalized theory of interference and its applications,''
Proc.\ Indian Acad.\ Sci.\ A \textbf{44}, 247--262 (1956).

\bibitem{Simon1983}
B.~Simon,
``Holonomy, the quantum adiabatic theorem, and Berry's phase,''
Phys.\ Rev.\ Lett.\ \textbf{51}, 2167--2170 (1983).

\bibitem{berry1984}
M.~V.~Berry,
``Quantal phase factors accompanying adiabatic changes,''
Proc.\ R.\ Soc.\ Lond.\ A \textbf{392}, 45--57 (1984).

\bibitem{WilczekZee1984}
F.~Wilczek and A.~Zee,
``Appearance of gauge structure in simple dynamical systems,''
Phys.\ Rev.\ Lett.\ \textbf{52}, 2111--2114 (1984).

\bibitem{AharonovAnandan1987}
Y.~Aharonov and J.~Anandan,
``Phase change during a cyclic quantum evolution,''
Phys.\ Rev.\ Lett.\ \textbf{58}, 1593--1596 (1987).

\bibitem{AnandanAharonov1990}
J.~Anandan and Y.~Aharonov,
``Geometry of quantum evolution,''
Phys.\ Rev.\ Lett.\ \textbf{65}, 1697--1700 (1990).

\bibitem{ProvostVallee1980}
J.~P.~Provost and G.~Vall\'ee,
``Riemannian structure on manifolds of quantum states,''
Commun.\ Math.\ Phys.\ \textbf{76}, 289--301 (1980).

\bibitem{Kahler1933}
E.~K\"ahler,
``\"Uber eine bemerkenswerte Hermitesche Metrik,''
Abh.\ Math.\ Sem.\ Univ.\ Hamburg \textbf{9}, 173--186 (1933).

\bibitem{Hopf1931}
H.~Hopf,
``\"Uber die Abbildungen der dreidimensionalen Sph\"are auf die Kugelfl\"ache,''
Math.\ Ann.\ \textbf{104}, 637--665 (1931).

\bibitem{Chern1946}
S.-S.~Chern,
``Characteristic classes of Hermitian manifolds,''
Ann.\ Math.\ \textbf{47}, 85--121 (1946).

\bibitem{AmbroseSinger1953}
W.~Ambrose and I.~M.~Singer,
``A theorem on holonomy,''
Trans.\ Am.\ Math.\ Soc.\ \textbf{75}, 428--443 (1953).

\bibitem{pontryagin1955}
L.~S.~Pontryagin,
``Smooth manifolds and their applications in homotopy theory,''
Trudy Mat.\ Inst.\ Steklov \textbf{45}, 1--139 (1955).

\bibitem{KobayashiNomizu1963}
S.~Kobayashi and K.~Nomizu,
\emph{Foundations of Differential Geometry}, Vols.~I--II
(Wiley-Interscience, New York, 1963--1969).

\bibitem{BottTu1982}
R.~Bott and L.~W.~Tu,
\emph{Differential Forms in Algebraic Topology}
(Springer, New York, 1982).

\bibitem{Nakahara2003}
M.~Nakahara,
\emph{Geometry, Topology and Physics}, 2nd ed.
(Institute of Physics Publishing, Bristol, 2003).

\bibitem{BengtssonZyczkowski2006}
I.~Bengtsson and K.~\.{Z}yczkowski,
\emph{Geometry of Quantum States}
(Cambridge University Press, Cambridge, 2006).

\item[]\textbf{E. Information theory, statistical geometry, and computation}

\bibitem{shannon1948}
C.~E.~Shannon,
``A mathematical theory of communication,''
Bell Syst.\ Tech.\ J.\ \textbf{27}, 379--423 and 623--656 (1948).

\bibitem{vonNeumann1932}
J.~von Neumann,
\emph{Mathematische Grundlagen der Quantenmechanik}
(Springer, Berlin, 1932).

\bibitem{Fisher1925}
R.~A.~Fisher,
``Theory of statistical estimation,''
Proc.\ Camb.\ Philos.\ Soc.\ \textbf{22}, 700--725 (1925).

\bibitem{Rao1945}
C.~R.~Rao,
``Information and the accuracy attainable in the estimation of statistical
parameters,'' Bull.\ Calcutta Math.\ Soc.\ \textbf{37}, 81--91 (1945).

\bibitem{Jaynes1957}
E.~T.~Jaynes,
``Information theory and statistical mechanics,''
Phys.\ Rev.\ \textbf{106}, 620--630 (1957); \textbf{108}, 171--190 (1957).

\bibitem{Landauer1961}
R.~Landauer,
``Irreversibility and heat generation in the computing process,''
IBM J.\ Res.\ Dev.\ \textbf{5}, 183--191 (1961).

\bibitem{Bures1969}
D.~Bures,
``An extension of Kakutani's theorem on infinite product measures to the
tensor product of semifinite $W^*$-algebras,''
Trans.\ Am.\ Math.\ Soc.\ \textbf{135}, 199--212 (1969).

\bibitem{Wootters1981}
W.~K.~Wootters,
``Statistical distance and Hilbert space,''
Phys.\ Rev.\ D \textbf{23}, 357--362 (1981).

\bibitem{Bennett1982}
C.~H.~Bennett,
``The thermodynamics of computation---a review,''
Int.\ J.\ Theor.\ Phys.\ \textbf{21}, 905--940 (1982).

\bibitem{BraunsteinCaves1994}
S.~L.~Braunstein and C.~M.~Caves,
``Statistical distance and the geometry of quantum states,''
Phys.\ Rev.\ Lett.\ \textbf{72}, 3439--3443 (1994).

\bibitem{AmariNagaoka2000}
S.-I.~Amari and H.~Nagaoka,
\emph{Methods of Information Geometry}
(American Mathematical Society, Providence, 2000).

\bibitem{NielsenChuang2000}
M.~A.~Nielsen and I.~L.~Chuang,
\emph{Quantum Computation and Quantum Information}
(Cambridge University Press, Cambridge, 2000).

\bibitem{Lloyd2000}
S.~Lloyd,
``Ultimate physical limits to computation,''
Nature \textbf{406}, 1047--1054 (2000).

\bibitem{CoverThomas2006}
T.~M.~Cover and J.~A.~Thomas,
\emph{Elements of Information Theory}, 2nd ed.
(Wiley, Hoboken, 2006).

\item[]\textbf{F. Synchronization, nonlinear dynamics, and collective phase models}

\bibitem{Kuramoto1975}
Y.~Kuramoto,
``Self-entrainment of a population of coupled nonlinear oscillators,'' in
\emph{International Symposium on Mathematical Problems in Theoretical Physics},
Lecture Notes in Physics Vol.~39 (Springer, Berlin, 1975), pp.~420--422.

\bibitem{Feigenbaum1978}
M.~J.~Feigenbaum,
``Quantitative universality for a class of nonlinear transformations,''
J.\ Stat.\ Phys.\ \textbf{19}, 25--52 (1978).

\bibitem{Strogatz2000}
S.~H.~Strogatz,
``From Kuramoto to Crawford: Exploring the onset of synchronization in
populations of coupled oscillators,''
Physica D \textbf{143}, 1--20 (2000).

\bibitem{Acebron2005}
J.~A.~Acebr\'on, L.~L.~Bonilla, C.~J.~P.~Vicente, F.~Ritort, and R.~Spigler,
``The Kuramoto model: A simple paradigm for synchronization phenomena,''
Rev.\ Mod.\ Phys.\ \textbf{77}, 137--185 (2005).

\bibitem{OttAntonsen2008}
E.~Ott and T.~M.~Antonsen,
``Low dimensional behavior of large systems of globally coupled oscillators,''
Chaos \textbf{18}, 037113 (2008).

\bibitem{Strogatz2015}
S.~H.~Strogatz,
\emph{Nonlinear Dynamics and Chaos}, 2nd ed.
(Westview Press, Boulder, 2015).

\item[]\textbf{G. Collatz dynamics, Ramanujan analysis, primes, and zeta structure}

\bibitem{Terras1976}
R.~Terras,
``A stopping time problem on the positive integers,''
Acta Arith.\ \textbf{30}, 241--252 (1976).

\bibitem{Lagarias1985}
J.~C.~Lagarias,
``The $3x+1$ problem and its generalizations,''
Am.\ Math.\ Mon.\ \textbf{92}, 3--23 (1985).

\bibitem{Wirsching1998}
G.~J.~Wirsching,
\emph{The Dynamical System Generated by the $3n+1$ Function}
(Springer, Berlin, 1998).

\bibitem{Tao2022Collatz}
T.~Tao,
``Almost all orbits of the Collatz map attain almost bounded values,''
Forum Math.\ Pi \textbf{10}, e12 (2022).

\bibitem{Ramanujan1918}
S.~Ramanujan,
``On certain trigonometrical sums and their applications in the theory of
numbers,'' Trans.\ Cambridge Philos.\ Soc.\ \textbf{22}, 259--276 (1918).

\bibitem{HardyRamanujan1918}
G.~H.~Hardy and S.~Ramanujan,
``Asymptotic formulae in combinatory analysis,''
Proc.\ Lond.\ Math.\ Soc.\ \textbf{17}, 75--115 (1918).

\bibitem{Rademacher1937}
H.~Rademacher,
``On the partition function $p(n)$,''
Proc.\ Natl.\ Acad.\ Sci.\ USA \textbf{23}, 78--84 (1937).

\bibitem{HardyLittlewood1923}
G.~H.~Hardy and J.~E.~Littlewood,
``Some problems of `Partitio Numerorum'; III: On the expression of a number as
a sum of primes,'' Acta Math.\ \textbf{44}, 1--70 (1923).

\bibitem{Riemann1859}
B.~Riemann,
``\"Uber die Anzahl der Primzahlen unter einer gegebenen Gr\"osse,''
Monatsber.\ K\"onigl.\ Preuss.\ Akad.\ Wiss.\ Berlin 671--680 (1859).

\bibitem{Selberg1956}
A.~Selberg,
``Harmonic analysis and discontinuous groups in weakly symmetric Riemannian
spaces with applications to Dirichlet series,''
J.\ Indian Math.\ Soc.\ \textbf{20}, 47--87 (1956).

\bibitem{BerryKeating1999}
M.~V.~Berry and J.~P.~Keating,
``The Riemann zeros and eigenvalue asymptotics,''
SIAM Rev.\ \textbf{41}, 236--266 (1999).

\bibitem{Connes1999}
A.~Connes,
``Trace formula in noncommutative geometry and the zeros of the Riemann zeta
function,'' Selecta Math.\ \textbf{5}, 29--106 (1999).

\bibitem{Zhang2014}
Y.~Zhang,
``Bounded gaps between primes,''
Ann.\ Math.\ \textbf{179}, 1121--1174 (2014).

\bibitem{Polymath2014}
D.~H.~J.~Polymath,
``New equidistribution estimates of Zhang type,''
Algebra Number Theory \textbf{8}, 2067--2199 (2014).

\bibitem{Maynard2015}
J.~Maynard,
``Small gaps between primes,''
Ann.\ Math.\ \textbf{181}, 383--413 (2015).

\bibitem{Apostol1976}
T.~M.~Apostol,
\emph{Introduction to Analytic Number Theory}
(Springer, New York, 1976).

\bibitem{HardyWright2008}
G.~H.~Hardy and E.~M.~Wright,
\emph{An Introduction to the Theory of Numbers}, 6th ed., revised by
D.~R.~Heath-Brown and J.~H.~Silverman
(Oxford University Press, Oxford, 2008).

\item[]\textbf{H. Anomalies, CP violation, strong CP, and electric dipole moments}

\bibitem{Adler1969}
S.~L.~Adler,
``Axial-vector vertex in spinor electrodynamics,''
Phys.\ Rev.\ \textbf{177}, 2426--2438 (1969).

\bibitem{BellJackiw1969}
J.~S.~Bell and R.~Jackiw,
``A PCAC puzzle: $\pi^0\to\gamma\gamma$ in the sigma model,''
Nuovo Cim.\ A \textbf{60}, 47--61 (1969).

\bibitem{tHooft1976}
G.~'t~Hooft,
``Symmetry breaking through Bell--Jackiw anomalies,''
Phys.\ Rev.\ Lett.\ \textbf{37}, 8--11 (1976).

\bibitem{PecceiQuinn1977}
R.~D.~Peccei and H.~R.~Quinn,
``CP conservation in the presence of pseudoparticles,''
Phys.\ Rev.\ Lett.\ \textbf{38}, 1440--1443 (1977).

\bibitem{Weinberg1978Axion}
S.~Weinberg,
``A new light boson?'' Phys.\ Rev.\ Lett.\ \textbf{40}, 223--226 (1978).

\bibitem{Wilczek1978Axion}
F.~Wilczek,
``Problem of strong $P$ and $T$ invariance in the presence of instantons,''
Phys.\ Rev.\ Lett.\ \textbf{40}, 279--282 (1978).

\bibitem{crewther1979}
R.~J.~Crewther, P.~Di~Vecchia, G.~Veneziano, and E.~Witten,
``Chiral estimate of the electric dipole moment of the neutron in quantum
chromodynamics,'' Phys.\ Lett.\ B \textbf{88}, 123--127 (1979).

\bibitem{witten1982}
E.~Witten,
``An SU(2) anomaly,'' Phys.\ Lett.\ B \textbf{117}, 324--328 (1982).

\bibitem{pospelov2005}
M.~Pospelov and A.~Ritz,
``Electric dipole moments as probes of new physics,''
Ann.\ Phys.\ (N.Y.) \textbf{318}, 119--169 (2005).

\item[]\textbf{I. Gravitation, cosmology, horizons, and information}

\bibitem{Einstein1915}
A.~Einstein,
``Die Feldgleichungen der Gravitation,''
Sitzungsber.\ Preuss.\ Akad.\ Wiss.\ Berlin 844--847 (1915).

\bibitem{Friedmann1922}
A.~Friedmann,
``\"Uber die Kr\"ummung des Raumes,''
Z.\ Phys.\ \textbf{10}, 377--386 (1922).

\bibitem{Lemaitre1927}
G.~Lema\^itre,
``Un univers homog\`ene de masse constante et de rayon croissant, rendant
compte de la vitesse radiale des n\'ebuleuses extra-galactiques,''
Ann.\ Soc.\ Sci.\ Bruxelles A \textbf{47}, 49--59 (1927).

\bibitem{Hubble1929}
E.~Hubble,
``A relation between distance and radial velocity among extra-galactic
nebulae,'' Proc.\ Natl.\ Acad.\ Sci.\ USA \textbf{15}, 168--173 (1929).

\bibitem{Penrose1965}
R.~Penrose,
``Gravitational collapse and space-time singularities,''
Phys.\ Rev.\ Lett.\ \textbf{14}, 57--59 (1965).

\bibitem{Bekenstein1973}
J.~D.~Bekenstein,
``Black holes and entropy,'' Phys.\ Rev.\ D \textbf{7}, 2333--2346 (1973).

\bibitem{Hawking1975}
S.~W.~Hawking,
``Particle creation by black holes,''
Commun.\ Math.\ Phys.\ \textbf{43}, 199--220 (1975).

\bibitem{HawkingEllis1973}
S.~W.~Hawking and G.~F.~R.~Ellis,
\emph{The Large Scale Structure of Space-Time}
(Cambridge University Press, Cambridge, 1973).

\bibitem{Wald1984}
R.~M.~Wald,
\emph{General Relativity}
(University of Chicago Press, Chicago, 1984).

\bibitem{Weinberg1989CC}
S.~Weinberg,
``The cosmological constant problem,''
Rev.\ Mod.\ Phys.\ \textbf{61}, 1--23 (1989).

\bibitem{Jacobson1995}
T.~Jacobson,
``Thermodynamics of spacetime: The Einstein equation of state,''
Phys.\ Rev.\ Lett.\ \textbf{75}, 1260--1263 (1995).

\bibitem{Carroll2001}
S.~M.~Carroll,
``The cosmological constant,''
Living Rev.\ Relativ.\ \textbf{4}, 1 (2001).

\bibitem{Verlinde2011}
E.~P.~Verlinde,
``On the origin of gravity and the laws of Newton,''
J.\ High Energy Phys.\ \textbf{04}, 029 (2011).

\item[]\textbf{J. Scientific methodology, falsifiability, preregistration, and reproducibility}

\bibitem{Popper1959}
K.~R.~Popper,
\emph{The Logic of Scientific Discovery}
(Hutchinson, London, 1959).

\bibitem{Lakatos1978}
I.~Lakatos,
\emph{The Methodology of Scientific Research Programmes}
(Cambridge University Press, Cambridge, 1978).

\bibitem{Ioannidis2005}
J.~P.~A.~Ioannidis,
``Why most published research findings are false,''
PLoS Med.\ \textbf{2}, e124 (2005).

\bibitem{Sandve2013}
G.~K.~Sandve, A.~Nekrutenko, J.~Taylor, and E.~Hovig,
``Ten simple rules for reproducible computational research,''
PLoS Comput.\ Biol.\ \textbf{9}, e1003285 (2013).

\bibitem{Wilson2014}
G.~Wilson \emph{et al.},
``Best practices for scientific computing,''
PLoS Biol.\ \textbf{12}, e1001745 (2014).

\bibitem{Munafo2017}
M.~R.~Munaf\`o \emph{et al.},
``A manifesto for reproducible science,''
Nat.\ Hum.\ Behav.\ \textbf{1}, 0021 (2017).

\bibitem{Nosek2018}
B.~A.~Nosek, C.~R.~Ebersole, A.~C.~DeHaven, and D.~T.~Mellor,
``The preregistration revolution,''
Proc.\ Natl.\ Acad.\ Sci.\ USA \textbf{115}, 2600--2606 (2018).

\bibitem{Stark2018}
P.~B.~Stark,
``Before reproducibility must come preproducibility,''
Nature \textbf{557}, 613 (2018).

\end{thebibliography}
